% from ega4-21c-fr.tex
\oldpage[IV-4]{324}
donc régulière, et en particulier intègre, autrement dit
$\mathfrak p=\mathfrak mB$ est bien un idéal premier de $B$, ce qui justifie
l'énoncé. Tout revient évidemment à prouver que tout idéal premier $\mathfrak q$ de
$B$ non contenu dans $\mathfrak p$ est \emph{principal}.

Soient $\widehat A$, $\widehat B$ les complétés de $A$ et $B$ respectivement, de
sorte que l'idéal maximal de $\widehat A$ est $\mathfrak m\widehat A$; on sait
$(\mathbf 0, 19.3.6)$ que $\widehat B$ est une $\widehat A$-algèbre formellement
lisse pour les topologies adiques. Soit $k'$ le corps résiduel de $\widehat B$,
extension finie de $k$; il existe un homomorphisme local $\widehat A\to C$, où $C$
est un anneau local noethérien qui est un $\widehat A$-module \emph{fini} et
\emph{plat} (donc libre) et est tel que $C/\mathfrak mC$ soit isomorphe à $k'$
$(\mathbf 0_{\mathrm{III}}, 10.3.1)$; on en déduit que $C$ est \emph{complet}, et il
résulte alors de (7.5.1), (7.5.3) et (6.5.4, (ii)) que $C$ est \emph{intègre et
intégralement clos}. En outre,
\[
  D=\widehat B\widehat\otimes_{\widehat A}C
\]
est un anneau semi-local complet, composé direct d'anneaux locaux complets dont l'un,
$D_0$, a pour corps résiduel $k'$ (puisque $k'\otimes_k k'$ est composé direct
d'anneaux locaux dont l'un est isomorphe à $k'$). Comme $D$ est formellement lisse
sur $C$, il en est de même de $D_0$; par suite $D_0/\mathfrak mD_0$ est une
$k'$-algèbre formellement lisse, de corps résiduel $k'$, ce qui entraîne qu'elle est
$k'$-isomorphe à une algèbre de séries formelles
$k'[[T_1,\ldots,T_n]]$ $(\mathbf 0, 19.6.4)$; on en conclut, par
$(\mathbf 0, 19.7.1.5)$, que $D_0$ est $C$-isomorphe à
$C[[T_1,\ldots,T_n]]$, et par suite \emph{intègre et intégralement clos} (Bourbaki,
\emph{Alg. comm.}, chap.~V, \S1, n\textsuperscript{o}~4, prop.~14). Comme les
morphismes $\operatorname{Spec}(D_0)\to\operatorname{Spec}(\widehat B)$ et
$\operatorname{Spec}(D_0)\to\operatorname{Spec}(B)$ sont fidèlement plats, on en
déduit que $B$ et $\widehat B$ sont aussi \emph{intègres et intégralement clos}
(2.1.13). Cela prouve que l'idéal $\mathfrak qD_0$ est divisoriel (Bourbaki,
\emph{Alg. comm.}, chap.~VII, \S1, n\textsuperscript{o}~10, prop.~15) et non contenu
dans $\mathfrak mD_0$, sans quoi on aurait
\[
  \mathfrak q=(\mathfrak qD_0)\cap B
  \subset(\mathfrak mD_0)\cap B=\mathfrak mB=\mathfrak p
\]
par fidèle platitude $(\mathbf 0_{\mathrm I}, 6.5.1)$. On en conclut qu'aucun des
idéaux premiers $\mathfrak r_h$ de hauteur $1$ dans $D_0$ qui contiennent
$\mathfrak qD_0$ ne peut être contenu dans $\mathfrak mD_0$; si l'on prouve que ces
idéaux sont principaux, il en résultera que $\mathfrak qD_0$ est principal, les
diviseurs (au sens de Bourbaki) de $\mathfrak qD_0$ et d'un produit de puissances des
$\mathfrak r_h$ étant égaux (Bourbaki, \emph{Alg. comm.}, chap.~VII, \S1,
n\textsuperscript{o}~4, prop.~5). Comme $\mathfrak qD_0\cap B=\mathfrak q$, on en
déduit par fidèle platitude que $\mathfrak q$ est principal (2.5.2), en utilisant le
fait que $B$ est un anneau local.

On est ainsi ramené à prouver le théorème dans le cas particulier où $A$ est
\emph{complet} et $B=A[[T_1,\ldots,T_n]]$. Montrons d'abord qu'on peut se ramener au
cas où $n=1$. En effet, procédons par récurrence sur $n$, et soit
$f\in\mathfrak q$ un élément n'appartenant pas à $\mathfrak p=\mathfrak mB$; il
suffira de montrer qu'il existe un $A$-automorphisme $\sigma$ de $B$ tel que
$\sigma(f)$ n'appartienne pas à l'idéal
\[
  \mathfrak p'=\mathfrak p+BT_1+\cdots+BT_{n-1};
\]
en effet, si $B'=A[[T_1,\ldots,T_{n-1}]]$, $B'$ est complet et a pour idéal maximal
$\mathfrak mB'+B'T_1+\cdots+B'T_{n-1}=\mathfrak m'$. Comme on a
$B=B'[[T_n]]$ et $\mathfrak p'=\mathfrak m'B$, il résultera de l'hypothèse de
récurrence (compte tenu de ce que $B'$ est intègre et intégralement clos) que l'idéal
$\sigma(\mathfrak q)$ est principal dans $B$, donc aussi $\mathfrak q$. Or, si
$\overline f\in k[[T_1,\ldots,T_n]]$ est l'image de $f$ (« série réduite » de $f$),
on sait (Bourbaki, \emph{Alg. comm.}, chap.~VII, \S3, n\textsuperscript{o}~7,
lemme~3) qu'il y a un $A$-automorphisme $\sigma$ de $B$ tel que
$(\overline{\sigma(f)})(0,\ldots,0,T_n)\neq0$, ce qui implique évidemment
$\sigma(f)\notin\mathfrak p'$.

Supposons donc $n=1$ et écrivons $T$ au lieu de $T_1$, de sorte que $B=A[[T]]$. Il
suffit de montrer que dans l'anneau de polynômes $A[T]$, l'idéal
$\mathfrak q\cap A[T]$ est principal;
